\pagestyle{empty}
\tight
\begin{tblr}{width=\textwidth,colspec={|Q[c,m,1.9cm]|X[l,m]|},hlines,vlines,hspan=minimal,row{1}={bg=headblue},rowsep=1pt,colsep=3pt}
\SetCell{c,m}\hfil\logo\hfil & \textbf{POLITEKNIK NEGERI MALANG}\par\textbf{JURUSAN TEKNOLOGI INFORMASI}\par\textbf{PROGRAM STUDI: D4 TEKNIK INFORMATIKA} \\
\SetCell[c=2]{c} \textbf{RENCANA PEMBELAJARAN SEMESTER (RPS)} & \\
\end{tblr}
\begin{longtblr}{width=\textwidth,colspec={|Q[l,m,1.9cm]|X[0.85,l,m]|X[1.45,l,m]|X[0.75,l,m]|X[1.15,l,m]|},hlines,vlines,hspan=minimal,rowsep=1pt,colsep=2pt,row{1,3}={bg=headgray,font=\bfseries}}
MATA KULIAH & KODE & BOBOT (SKS)/jam & SEMESTER & TGL. PENYUSUNAN \\
Basis Data Lanjut & RTI253005 & 2 SKS/4 jam & 3 & 6 Juli 2026 \\
OTORISASI & Dosen Pengembang RPS & Koordinator RMK & \SetCell[c=2]{l} Ka PRODI &  \\
 & Candra Bella Vista, S.Kom., M.T. & Candra Bella Vista, S.Kom., M.T. & \SetCell[c=2]{l}{\vspace{8mm}\par Dr. Ely Setyo Astuti, S.T., M.T.} &  \\
\SetCell[r=12]{l} Capaian Pembelajaran (CP) & \SetCell[c=4]{l,bg=headgray,font=\bfseries} Capaian Pembelajaran Lulusan Program Studi (CPL-Prodi) &  &  &  \\
 & CPL02 & \SetCell[c=3]{l} Mampu menerapkan prinsip rekayasa sistem dan perangkat lunak untuk menghasilkan solusi aplikasi multi-platform sesuai kebutuhan. &  &  \\
 & CPL07 & \SetCell[c=3]{l} Mampu membangun sistem komputasi cerdas dan menganalisis data berbasis kecerdasan buatan untuk mendukung pengambilan keputusan dan inovasi. &  &  \\
 & CPL10 & \SetCell[c=3]{l} Memiliki kompetensi untuk menganalisis persoalan kompleks untuk mengidentifikasi solusi bidang informatika/ilmu komputer dengan mempertimbangkan wawasan ilmu multidisiplin. &  &  \\
 & \SetCell[c=4]{l,bg=headgray,font=\bfseries} Capaian Pembelajaran mata kuliah (CPL-MK) &  &  &  \\
 & CPMK0203 & \SetCell[c=3]{l} Mampu merancang skema dan arsitektur manajemen data yang efektif dan terstruktur untuk mendukung solusi aplikasi multi-platform dengan mempertimbangkan integritas data, kebutuhan informasi stakeholder, serta efisiensi dalam operasi query dan skalabilitas sistem. &  &  \\
 & CPMK0701 & \SetCell[c=3]{l} Mampu membangun sistem manajemen data yang handal untuk penyimpanan dan pengelolaan informasi yang efisien, termasuk optimasi performa dan integritas data. &  &  \\
 & CPMK1001 & \SetCell[c=3]{l} Mampu menganalisis permasalahan pengelolaan data dan informasi yang bersifat kompleks dengan wawasan multidisiplin untuk merancang serta mengimplementasikan solusi berbasis basis data yang sesuai dengan kebutuhan organisasi atau bisnis. &  &  \\
 & \SetCell[c=4]{l,bg=headgray,font=\bfseries} Kemampuan akhir tiap tahapan belajar (Sub-CPMK) &  &  &  \\
 & SCPMK0203-02301 & \SetCell[c=3]{l} Mampu mengoperasikan PostgreSQL dan membuat query dasar (DDL, DML) serta query lanjutan (sub-query, JOIN, grouping, dan CTE) untuk manajemen data. &  &  \\
 & SCPMK0203-02302 & \SetCell[c=3]{l} Mampu mengimplementasikan teknik optimasi query menggunakan indeks serta membangun objek basis data (function, view, materialized view, dan stored procedure) pada PostgreSQL. &  &  \\
 & SCPMK0701-02303 & \SetCell[c=3]{l} Mampu menerapkan mekanisme transaksi, concurrency, full-text search, JSONB, backup-restore, dan migrasi basis data pada PostgreSQL. &  &  \\
 & SCPMK1001-02304 & \SetCell[c=3]{l} Mampu mengintegrasikan PostgreSQL dengan aplikasi dan menyelesaikan studi kasus kompleks melalui proyek basis data terintegrasi. &  &  \\
\SetCell[r=2]{l} Penilaian & \SetCell[c=4]{l} *Nilai CPMK didapat dari agregasi nilai SUB-CPMK yang dibebankan &  &  &  \\
 & \SetCell[c=4]{l}{\tiny\begin{tblr}{width=\linewidth,colspec={|Q[c,m,0.1\linewidth]|Q[c,m,0.16\linewidth]|X[l,m]|X[l,m]|X[l,m]|X[l,m]|X[l,m]|Q[c,m,0.08\linewidth]|},hlines,vlines,hspan=minimal,rowsep=1pt,colsep=0.7pt,row{1,2}={bg=headgray,font=\bfseries}}
Kode CPMK & Kode SCPMK & \SetCell[c=5]{c} Bobot per Bentuk Penilaian & & & & & Total Bobot \\
 & & Tugas & Kuis 1 & UTS & PBL & UAS & \\
\SetCell[r=2]{c} CPMK0203 & SCPMK0203-02301 & 10\%: jobsheet query dasar dan lanjutan & 10\%: konsep PostgreSQL dan SQL & 5\%: query dasar dan CTE & & & 25\% \\
 & SCPMK0203-02302 & 5\%: jobsheet index dan objek basis data & & 8\%: index dan stored procedure & & & 13\% \\
\SetCell[r=1]{c} CPMK0701 & SCPMK0701-02303 & 5\%: jobsheet transaksi dan fitur lanjutan & & 7\%: transaksi dan JSONB & & & 12\% \\
\SetCell[r=1]{c} CPMK1001 & SCPMK1001-02304 & 5\%: jobsheet studi kasus integrasi & & & 20\%: proyek basis data terintegrasi & 25\%: implementasi database PBL & 50\% \\
\SetCell[c=7]{r}\textbf{Total Per Penilaian} & & & & & & & \textbf{100\%} \\
\end{tblr}} &  &  &  \\
Diskripsi Singkat Mata Kuliah & \SetCell[c=4]{l} Mata kuliah Basis Data Lanjut membangun kompetensi mahasiswa dalam mengelola dan mengoptimalkan basis data menggunakan PostgreSQL. Topik meliputi query lanjutan (sub-query, JOIN, CTE), optimasi dengan indeks, pembuatan function, view, stored procedure, manajemen transaksi dan concurrency, full-text search, JSONB, backup-restore, migrasi dari MySQL, serta integrasi basis data dengan aplikasi dalam studi kasus nyata. &  &  &  \\
Bahan Kajian / Materi Pembelajaran & \SetCell[c=4]{l}\begin{enumerate}\item Pengenalan PostgreSQL dan query dasar\item Query lanjutan: sub-query, JOIN, grouping, dan CTE\item Indeks dan optimasi query\item Function, view, materialized view, dan stored procedure\item Transaksi, concurrency, full-text search, dan JSONB\item Backup-restore dan migrasi basis data\item Integrasi basis data dengan aplikasi dan studi kasus\end{enumerate} &  &  &  \\
\SetCell[r=2]{l} Pustaka & \SetCell[c=4]{l} \textbf{Utama:}\begin{enumerate}\item Ramakrishnan, R. \& Gehrke, J. 2003. \textit{Database Management Systems}, 3rd ed. McGraw-Hill.\item PostgreSQL Global Development Group. 2024. \textit{PostgreSQL 16 Documentation}. https://www.postgresql.org/docs/.\item Elmasri, R. \& Navathe, S. B. 2016. \textit{Fundamentals of Database Systems}, 7th ed. Pearson.\end{enumerate} &  &  &  \\
 & \SetCell[c=4]{l} \textbf{Pendukung:}\begin{enumerate}\item Silberschatz, A., Korth, H., Sudarshan, S. 2019. \textit{Database System Concepts}, 7th ed. McGraw-Hill.\item Momjian, B. 2001. \textit{PostgreSQL: Introduction and Concepts}. Addison-Wesley.\item Materi LMS, dokumentasi resmi PostgreSQL, dan jobsheet praktikum.\end{enumerate} &  &  &  \\
Media Pembelajaran & \SetCell[c=2]{l} \textbf{Software:} PostgreSQL, pgAdmin/DBeaver, VS Code, LMS. &  & \SetCell[c=2]{l} \textbf{Hardware:} Komputer/laptop, proyektor. &  \\
Nama Dosen Pengampu & \SetCell[c=4]{l}\begin{enumerate}\item Candra Bella Vista, S.Kom., M.T.\item M. Hasyim Ratsanjani, S.Kom., M.Kom.\item Ridwan Rismanto, S.ST., M.Kom., Ph.D.\end{enumerate} &  &  &  \\
Matakuliah Syarat & \SetCell[c=4]{l} RTI252006 Basis Data &  &  &  \\
\end{longtblr}
\begin{longtblr}{width=\textwidth,colspec={|Q[c,m,0.06\textwidth]|X[1.35,l,m]|X[1.08,l,m]|X[1.16,l,m]|X[0.7,l,m]|X[1.24,l,m]|X[1.06,l,m]|X[0.98,l,m]|Q[c,m,0.05\textwidth]|},hlines,vlines,hspan=minimal,rowhead=2,row{1,2}={bg=headgreen,font=\bfseries},rowsep=1pt,colsep=1.5pt}
Mgg.\par Ke & Kemampuan Akhir Yang Direncanakan (Sub-CP-MK) & Bahan kajian (Materi Pembelajaran) & Bentuk dan Metode Pembelajaran & Estimasi Waktu & Pengalaman Belajar Mahasiswa & Kriteria \& Bentuk Penilaian & Indikator Penilaian & Bobot Penilaian (\%) \\
(1) & (2) & (3) & (4) & (5) & (6) & (7) & (8) & (9) \\
1 & SCPMK0203-02301 & Pengenalan PostgreSQL: arsitektur, instalasi, konfigurasi awal, psql, pgAdmin, dan perbedaan dengan MySQL. & Modalitas: Blended Learning. Bentuk: Luring/Daring. Metode: Problem-Based Learning, praktikum laboratorium, demonstrasi, dan peer review. & 1 x 2 x 50' tatap muka; 1 x 2 x 50' tugas/praktik mandiri. & Mahasiswa menginstalasi PostgreSQL, menjelajahi antarmuka psql dan pgAdmin, serta mendokumentasikan perbedaan arsitektur dengan MySQL. & Bentuk penilaian: Jobsheet pengenalan PostgreSQL. Mengacu rubrik RTM. & Ketepatan instalasi dan eksplorasi antarmuka; dokumentasi perbedaan DBMS. & 2.5 \\
2 & SCPMK0203-02301 & Query Dasar DDL dan DML: CREATE, ALTER, DROP, INSERT, UPDATE, DELETE, SELECT dasar dengan klausa WHERE, ORDER BY, GROUP BY. & Modalitas: Blended Learning. Bentuk: Luring/Daring. Metode: Problem-Based Learning, praktikum laboratorium, demonstrasi, dan peer review. & 1 x 2 x 50' tatap muka; 1 x 2 x 50' tugas/praktik mandiri. & Mahasiswa membuat skema basis data, mengelola data menggunakan DDL dan DML, serta mengeksekusi query SELECT dasar berdasarkan studi kasus. & Bentuk penilaian: Jobsheet DDL dan DML. Mengacu rubrik RTM. & Ketepatan sintaks DDL/DML; kelengkapan implementasi skema dan data. & 2.5 \\
3 & SCPMK0203-02301 & Query Lanjutan: sub-query skalar dan korelasi, INNER/LEFT/RIGHT/FULL OUTER JOIN, CROSS JOIN, GROUP BY dengan HAVING, dan Common Table Expression (CTE). & Modalitas: Blended Learning. Bentuk: Luring/Daring. Metode: Problem-Based Learning, praktikum laboratorium, demonstrasi, dan peer review. & 1 x 2 x 50' tatap muka; 1 x 2 x 50' tugas/praktik mandiri. & Mahasiswa menyusun query lanjutan menggunakan sub-query, berbagai jenis JOIN, dan CTE untuk mengambil data dari beberapa tabel sesuai studi kasus. & Bentuk penilaian: Jobsheet query lanjutan. Mengacu rubrik RTM. & Ketepatan query sub-query dan JOIN; penggunaan CTE yang tepat. & 5 \\
4 & SCPMK0203-02301 & Kuis 1: materi minggu 1--3 (pengenalan PostgreSQL, DDL, DML, sub-query, JOIN, CTE). & Modalitas: Luring. Bentuk: Ujian tertulis/praktik, sifat closed book. & Sesuai jadwal asesmen. & Mahasiswa mengerjakan soal kuis berbasis kasus SQL menggunakan PostgreSQL secara mandiri. & Bentuk penilaian: Kuis 1 SQL dasar dan lanjutan. Mengacu rubrik RTM. & Ketepatan jawaban SQL; kemampuan menerapkan query pada kasus. & 10 \\
5 & SCPMK0203-02302 & Indeks dan Optimasi Query: jenis indeks (B-tree, Hash, GiST, GIN), EXPLAIN ANALYZE, query plan, identifikasi bottleneck, dan strategi optimasi. & Modalitas: Blended Learning. Bentuk: Luring/Daring. Metode: Problem-Based Learning, praktikum laboratorium, demonstrasi, dan peer review. & 1 x 2 x 50' tatap muka; 1 x 2 x 50' tugas/praktik mandiri. & Mahasiswa membuat berbagai jenis indeks, menganalisis query plan menggunakan EXPLAIN ANALYZE, dan mengidentifikasi serta mengoptimalkan query yang lambat. & Bentuk penilaian: Jobsheet indeks dan optimasi. Mengacu rubrik RTM. & Ketepatan pembuatan indeks; analisis query plan dan hasil optimasi. & 2.5 \\
6 & SCPMK0203-02302 & Function, View, Materialized View, dan Stored Procedure: pembuatan PL/pgSQL function, view, materialized view (REFRESH), dan stored procedure untuk enkapsulasi logika bisnis. & Modalitas: Blended Learning. Bentuk: Luring/Daring. Metode: Problem-Based Learning, praktikum laboratorium, demonstrasi, dan peer review. & 1 x 2 x 50' tatap muka; 1 x 2 x 50' tugas/praktik mandiri. & Mahasiswa membuat function, view, materialized view, dan stored procedure menggunakan PL/pgSQL sesuai studi kasus, lalu menguji dan mendokumentasikan hasilnya. & Bentuk penilaian: Jobsheet objek basis data. Mengacu rubrik RTM. & Ketepatan pembuatan function dan stored procedure; kebenaran view dan refresh materialized view. & 2.5 \\
7 & SCPMK0701-02303 & Transaksi dan Concurrency: ACID properties, BEGIN/COMMIT/ROLLBACK, SAVEPOINT, isolation level, serta penanganan deadlock dan race condition di PostgreSQL. & Modalitas: Blended Learning. Bentuk: Luring/Daring. Metode: Problem-Based Learning, praktikum laboratorium, demonstrasi, dan peer review. & 1 x 2 x 50' tatap muka; 1 x 2 x 50' tugas/praktik mandiri. & Mahasiswa mengimplementasikan transaksi ACID, menguji berbagai isolation level, serta mensimulasikan dan menangani deadlock pada skenario concurrency. & Bentuk penilaian: Jobsheet transaksi dan concurrency. Mengacu rubrik RTM. & Ketepatan implementasi transaksi ACID; penanganan isolation level dan deadlock. & 2.5 \\
8 & SCPMK0701-02303 & Full-Text Search dan JSONB: tsvector, tsquery, GIN index untuk FTS, operator JSONB (containment, path expression), indexing JSONB, dan query data semi-terstruktur. & Modalitas: Blended Learning. Bentuk: Luring/Daring. Metode: Problem-Based Learning, praktikum laboratorium, demonstrasi, dan peer review. & 1 x 2 x 50' tatap muka; 1 x 2 x 50' tugas/praktik mandiri. & Mahasiswa mengimplementasikan full-text search menggunakan tsvector/tsquery dan memanipulasi data JSONB menggunakan operator dan path expression pada PostgreSQL. & Bentuk penilaian: Jobsheet FTS dan JSONB. Mengacu rubrik RTM. & Ketepatan implementasi full-text search; kebenaran query JSONB dan indexing. & 2.5 \\
9 & SCPMK0203-02301, SCPMK0203-02302, SCPMK0701-02303 & UTS: materi minggu 1--8 (query dasar, query lanjutan, indeks, objek basis data, transaksi, FTS, JSONB). & Modalitas: Luring. Bentuk: Ujian praktik PostgreSQL, sifat closed book. & Sesuai jadwal asesmen. & Mahasiswa mengerjakan soal UTS berbasis studi kasus menggunakan PostgreSQL secara mandiri. & Bentuk penilaian: UTS query PostgreSQL lanjutan. Mengacu rubrik RTM. & Ketepatan query; kualitas optimasi dan implementasi objek basis data. & 20 \\
10 & SCPMK1001-02304 & Backup-Restore dan Migrasi MySQL ke PostgreSQL: pg\_dump, pg\_restore, psql dump, strategi backup, pgloader, konversi tipe data, dan penanganan perbedaan sintaks. & Modalitas: Blended Learning. Bentuk: Luring/Daring. Metode: Problem-Based Learning, praktikum laboratorium, demonstrasi, dan peer review. & 1 x 2 x 50' tatap muka; 1 x 2 x 50' tugas/praktik mandiri. & Mahasiswa melakukan backup dan restore basis data PostgreSQL, kemudian memigrasikan basis data dari MySQL menggunakan pgloader dan memverifikasi hasilnya. & Bentuk penilaian: Jobsheet backup-restore dan migrasi. Mengacu rubrik RTM. & Ketepatan proses backup-restore; keberhasilan migrasi dan validasi data. & 1.25 \\
11 & SCPMK1001-02304 & Studi Kasus Integrasi dengan Aplikasi: koneksi PostgreSQL dari Python/Node.js/PHP, ORM dasar (SQLAlchemy/Prisma), connection pool, dan penanganan error query. & Modalitas: Blended Learning. Bentuk: Luring/Daring. Metode: Problem-Based Learning, praktikum laboratorium, demonstrasi, dan peer review. & 1 x 2 x 50' tatap muka; 1 x 2 x 50' tugas/praktik mandiri. & Mahasiswa mengintegrasikan PostgreSQL dengan aplikasi menggunakan bahasa pemrograman pilihan, mengimplementasikan operasi CRUD, dan mendokumentasikan kode beserta hasilnya. & Bentuk penilaian: Jobsheet integrasi aplikasi. Mengacu rubrik RTM. & Ketepatan integrasi aplikasi-database; kelengkapan operasi CRUD dan dokumentasi. & 1.25 \\
12 & SCPMK1001-02304 & Studi Kasus 1: analisis kebutuhan dan perancangan skema basis data PostgreSQL untuk domain e-commerce. & Modalitas: Blended Learning. Bentuk: Luring/Daring. Metode: Problem-Based Learning, kolaborasi kelompok, dan presentasi. & 1 x 2 x 50' tatap muka; 1 x 2 x 50' tugas/praktik mandiri. & Mahasiswa menganalisis kebutuhan data domain e-commerce, merancang skema basis data PostgreSQL dengan optimasi indeks, lalu mengimplementasikan dan menguji query yang diperlukan. & Bentuk penilaian: Studi kasus 1 (bagian PBL). Mengacu rubrik RTM. & Ketepatan rancangan skema; kualitas implementasi dan query. & 1.25 \\
13 & SCPMK1001-02304 & Studi Kasus 2: implementasi stored procedure, trigger, dan function untuk domain manajemen inventori. & Modalitas: Blended Learning. Bentuk: Luring/Daring. Metode: Problem-Based Learning, kolaborasi kelompok, dan presentasi. & 1 x 2 x 50' tatap muka; 1 x 2 x 50' tugas/praktik mandiri. & Mahasiswa mengimplementasikan stored procedure, trigger, dan function untuk otomasi proses bisnis pada domain manajemen inventori, kemudian menguji dan mendokumentasikan hasilnya. & Bentuk penilaian: Studi kasus 2 (bagian PBL). Mengacu rubrik RTM. & Ketepatan logika stored procedure dan trigger; pengujian dan dokumentasi. & 1.25 \\
14 & SCPMK1001-02304 & Studi Kasus 3: optimasi performa basis data menggunakan indeks, query tuning, dan materialized view untuk domain analitik. & Modalitas: Blended Learning. Bentuk: Luring/Daring. Metode: Problem-Based Learning, kolaborasi kelompok, dan presentasi. & 1 x 2 x 50' tatap muka; 1 x 2 x 50' tugas/praktik mandiri. & Mahasiswa menganalisis bottleneck performa, mengimplementasikan strategi optimasi (indeks, query tuning, materialized view), dan mengukur peningkatan performa. & Bentuk penilaian: Studi kasus 3 (bagian PBL). Mengacu rubrik RTM. & Ketepatan analisis dan implementasi optimasi; pengukuran peningkatan performa. & 1.25 \\
15 & SCPMK1001-02304 & Studi Kasus 4: implementasi full-text search dan pengelolaan data JSONB untuk domain portal berita. & Modalitas: Blended Learning. Bentuk: Luring/Daring. Metode: Problem-Based Learning, kolaborasi kelompok, dan presentasi. & 1 x 2 x 50' tatap muka; 1 x 2 x 50' tugas/praktik mandiri. & Mahasiswa mengimplementasikan full-text search dan pengelolaan data semi-terstruktur JSONB pada domain portal berita, termasuk pembuatan indeks GIN dan query kompleks. & Bentuk penilaian: Studi kasus 4 (bagian PBL). Mengacu rubrik RTM. & Ketepatan implementasi FTS dan JSONB; kualitas query dan indeks. & 1.25 \\
16 & SCPMK1001-02304 & Studi Kasus 5: integrasi penuh basis data PostgreSQL dengan aplikasi web dalam skenario multi-pengguna dengan manajemen transaksi. & Modalitas: Blended Learning. Bentuk: Luring/Daring. Metode: Problem-Based Learning, kolaborasi kelompok, dan presentasi. & 1 x 2 x 50' tatap muka; 1 x 2 x 50' tugas/praktik mandiri. & Mahasiswa mengintegrasikan semua fitur PostgreSQL dalam aplikasi web multi-pengguna dengan manajemen transaksi yang tepat, kemudian mempresentasikan rancangan dan hasilnya. & Bentuk penilaian: Studi kasus 5 (bagian PBL). Mengacu rubrik RTM. & Ketepatan integrasi dan manajemen transaksi; kualitas presentasi. & 1.25 \\
17 & SCPMK1001-02304 & Proyek Akhir (UAS): implementasi proyek basis data PostgreSQL terintegrasi sesuai topik PBL kelompok, mencakup perancangan, implementasi, optimasi, dan presentasi. & Modalitas: Luring. Bentuk: Ujian praktik dan defense, sifat open problem. & Sesuai jadwal asesmen. & Mahasiswa mempresentasikan dan mendemonstrasikan proyek basis data PostgreSQL terintegrasi yang mencakup seluruh kompetensi mata kuliah. & Bentuk penilaian: UAS proyek akhir. Mengacu rubrik RTM. & Kelengkapan fitur proyek; kualitas implementasi dan kemampuan defense. & 25 \\
\end{longtblr}
